% SGA 3, Expose VI_A, Loop 1 English reconstruction.
% Scope: Examples 1.3.2 through the end of Corollary 2.4.1;
% paragraph 2.5 is excluded.
% Authority: Polo--Gille Exp6A-13oct24.pdf, component-local pages 7--11
% (combined-reader pages 310--314; printed pages 299--303).
% The supplied opening locator (local p. 6 / combined p. 309 / printed
% p. 298) is the immediately preceding page; the Examples 1.3.2 heading
% is visibly printed on the next page.
% OCR and the Reinhold English Markdown are comparison witnesses only.

\medskip
\phantomsection\label{ex:sga3-VIA-1.3.2}%
\noindent\textbf{例 1.3.2.}
设 \(k\) 是特征为 \(p>0\) 的非完美域,
\(t\in k-k^p\), \(\overline k\) 是 \(k\) 的一个代数闭包,
且 \(\alpha\in\overline k\) 满足 \(\alpha^p=t\).

\begin{enumerate}
\item[(1)] 考虑加法群
\(\mathbb G_{a,k}=\operatorname{Spec}k[X]\),
设 \(G\) 是由加法多项式 \(X^{p^2}-tX^p\) 定义的群子概形,
它在 \(k\) 上有限. 则
\[
  G_{\mathrm{red}}
    =\operatorname{Spec}k[X]/X\bigl(X^{p(p-1)}-t\bigr)
\]
在原点处平展. 若它是 \(k\) 上的群概形,
则它将平滑 (由 \hyperref[prop:sga3-VIA-1.3.1]{1.3.1});
但 \(G_{\mathrm{red}}\) 并非几何约化,
所以它不是 \(k\) 上的群概形.

\item[(2)] 考虑
\[
  \mathbb G^4_{a,k}=\operatorname{Spec}k[X,Y,U,V],
\]
设 \(G\) 是由加法多项式
\[
  P=X^p-tY^p
  \qquad\text{及}\qquad
  Q=U^p-tV^p.
\]
所生成的理想 \(I\) 定义的群子概形.
则 \(G\) 是二维且不可约的, 因为
\[
  (G_{\overline k})_{\mathrm{red}}
    \simeq\operatorname{Spec}\overline k[Y,V]
\]
不可约.

设 \(A=k[X,Y,U,V]\), \(\mathfrak m\) 是其增广理想.
以 \(x,y,u,v\) 表示 \(dX,dY,dU,dV\) 在
\[
  \Omega^1_{A/k}\otimes_A(A/\mathfrak m),
\]
中的像, 将它们视为切空间
\[
  k^4=T_0\mathbb G^4_{a,k}.
\]
上的线性形式.
我们要证明子空间 \(E=T_0G_{\mathrm{red}}\) 等于 \(k^4\).
否则, 将有一个系数不全为零的线性形式
\[
  f=ax+by+a'u+b'v,
  \qquad a,b,a',b'\in k,
\]
在 \(E\) 上为零.
回顾 \(\Omega^1_{A/k}\) 的构造, 从而切空间的构造,
与基变换可交换 (参见 EGA~IV\textsubscript{4}, 16.4.5),
并将 \(f\) 等同于其在 \((\overline k^4)^*\) 中的像.
由于
\[
  (G_{\overline k})_{\mathrm{red}}
    \subset (G_{\mathrm{red}})_{\overline k},
\]
线性形式 \(f\) 在 \(\overline k^4\) 的子空间
\(T_0(G_{\overline k})_{\mathrm{red}}\) 上为零,
而该子空间由方程
\[
  g_1=x-\alpha y
  \qquad\text{及}\qquad
  g_2=u-\alpha v.
\]
定义. 因此
\[
  f=\lambda g_1+\mu g_2,
  \qquad \lambda,\mu\in\overline k.
\]
但 \(\lambda g_1+\mu g_2\) 属于 \(k^4\),
仅当 \(\lambda=\mu=0\).
此矛盾表明 \(E=k^4\), 因而
\[
  T_0(G_{\mathrm{red}})_{\overline k}=\overline k^4.
\]

另一方面, \(R=XV-YU\) 属于 \(\sqrt I\), 因为
\[
  R^p=(X^p-tY^p)V^p-Y^p(U^p-tV^p).
\]
所以, \((G_{\mathrm{red}})_{\overline k}\) 在点
\((\alpha,1,\alpha,1)\) 处的切空间 \(F\),
包含在 \(\overline k^4\) 中方程为
\[
  \alpha\,dV+dX-dU-\alpha\,dY=0,
\]
的超平面 \(H\) 内, 因而维数至多为 \(3\).%
\footnote{\label{note:sga3-VIA-16}编者注: 事实上, 容易看出,
在每个非 \(0\) 的点处, \(\sqrt I\) 都由 \(P,Q,R\) 生成,
所以 \(F=H\).}
由 \SGARefLink{prop:sga3-VIA-1.3.1}{1.3.1} 的
\hyperref[concl:sga3-VIA-1.3.1-1]{第 (1) 部分},
可知 \(G_{\mathrm{red}}\) 不是 \(k\) 上的群概形.
\end{enumerate}

\section*{2. 局部有限型 \(A\)-群的连通分支}
\phantomsection\label{sec:sga3-VIA-2}
\addcontentsline{toc}{section}{2. 局部有限型
  \texorpdfstring{\(A\)}{A}-群的连通分支}

\medskip
\phantomsection\label{par:sga3-VIA-2.1}%
\noindent\textbf{2.1.}
先考虑任意 \(A\)-群 \(G\), 以 \(G'\) 表示 \(G\)
中原点 \(e\) 所在的连通分支.
此连通分支显然是闭的,
所以可以将它等同于以 \(G'\) 为底空间的 \(G\) 的约化闭子概形.%
\footnote{\label{note:sga3-VIA-17}编者注:
下文 (\SGARefLink{par:sga3-VIA-2.2}{2.2}) 将说明,
\(G'\) 是范畴 \((\mathbf{Sch}_{/k})_{\mathrm{red}}\) 中的一个群.}

\medskip
\phantomsection\label{prop:sga3-VIA-2.1.1}%
\noindent\textbf{命题 2.1.1.}
\textit{对 \(A\) 的剩余域 \(k\) 的每个扩张 \(K\),
\(G'\otimes_AK\) 的底空间是 \(K\)-群 \(G\otimes_AK\)
中原点所在的连通分支; 换言之, \(G'\) 几何连通.}

\medskip
事实上, 以 \((G\otimes_AK)'\) 表示 \(G\otimes_AK\)
中原点所在的连通分支.
由于 \((G\otimes_AK)'\) 在 \(G\) 中的像连通,
且包含 \(G\) 的单位元, 此像包含在 \(G'\) 中.
所以, \((G\otimes_AK)'\) 包含在 \(G'\) 于
\(G\otimes_AK\) 中的逆像 \(G'\otimes_AK\) 中.
因此, 只要证明 \(G'\otimes_AK\) 连通便得到此命题,
而这由下列引理证明.

\medskip
\phantomsection\label{lem:sga3-VIA-2.1.2}%
\noindent\textbf{引理 2.1.2.}
\textit{设 \(X\) 和 \(Y\) 是域 \(k\) 上的两个连通概形.
若 \(X\) 包含一个有理点, 则 \(X\times_kY\) 连通.}

\medskip
以下直接证明 EGA~IV\textsubscript{2} (4.5.8 和 4.5.14) 中的这一结果.

先假设 \(Y\) 非空、连通且仿射, 其代数为 \(B\).
则 \(X\times_kY\) 是拟凝聚 \(\mathcal O_X\)-代数
\[
  \mathcal B=\mathcal O_X\otimes_kB.
\]
的谱. 我们要证明, \(X\times_kY\) 的每个非空既开又闭子集
\(U\) 都等于 \(X\times_kY\).
现在, \(U\) 在 \(X\) 上仿射,
其仿射 \(\mathcal O_X\)-代数是 \(\mathcal B\) 的一个直因子.
由下文\hyperref[lem:sga3-VIA-2.1.3]{引理 2.1.3},
\(U\) 在 \(X\) 中的像既开又闭, 所以就是整个 \(X\).
特别地, 这个像包含 \(X\) 的一个有理点 \(x\),
所以 \(U\) 与 \(x\) 在 \(X\times_kY\) 中的逆像相交.
由于此逆像同构于 \(Y\), 因而连通,
\(U\) 包含整个逆像.
若 \(U\) 不同于 \(X\times_kY\),
对 \(U\) 在 \(X\times_kY\) 中的补集也有同样结论, 矛盾.

现在若 \(Y\) 是任意 \(k\)-概形,
以上论证表明典范投影
\[
  X\times_kY\longrightarrow Y
\]
的纤维连通. 若 \(x\) 是 \(X\) 的一个有理点,
则所有这些纤维都与子概形 \(\{x\}\times_kY\) 相交,
而该子概形本身连通. 断言得证.

\medskip
\phantomsection\label{lem:sga3-VIA-2.1.3}%
\noindent\textbf{引理 2.1.3.}
\textit{设 \(X\) 是一个概形, \(\mathcal A\) 是一个拟凝聚
\(\mathcal O_X\)-代数, 它局部地%
\footnote{\label{note:sga3-VIA-18}编者注: ``局部地'' 一词为新增内容.
对 Kaplansky 的引用是 \emph{Projective modules},
Ann.\ of Math.\ \textbf{68} (1958), 372--377;
另见 \textup{[BAC]}, \S~II.3, 习题 3.}
是一个自由 \(\mathcal O_X\)-模的直因子.
则 \(\operatorname{Spec}\mathcal A\) 在 \(X\) 中的像既开又闭.}

\medskip
设 \(V\) 为此像.
显然 \(V\) 包含在 \(\mathcal A\) 的支集中.
反之, 若 \(x\) 属于 \(\mathcal A\) 的支集,
则 \(\mathcal A_x\) 非零, 且为自由 \(\mathcal O_{X,x}\)-模,
因为由 Kaplansky 的结果, 局部环上的每个投射模都是自由的.
所以, \(\operatorname{Spec}\mathcal A\) 在 \(x\) 处的纤维是非空的,
它是代数为
\[
  \mathcal A_x\otimes_{\mathcal O_{X,x}}\kappa(x),
\]
的仿射概形.
因此, \(\operatorname{Spec}\mathcal A\) 的像与 \(\mathcal A\) 的支集重合.

若 \(s\) 是 \(\mathcal A\) 的单位截面,
则等式 \(s_x=0\) 蕴含 \(s\), 从而 \(\mathcal A\),
在 \(x\) 的一个邻域中为零.
因此, \(\mathcal A\) 的支集是闭的.
另一方面, 可以假设 \(\mathcal A\) 是 \(\mathcal L\) 的一个直因子,
后者是由 \(\mathcal O_X\) 的副本组成的族
\((\mathcal O_X^\alpha)\) 的直和.%
\footnote{\label{note:sga3-VIA-19}编者注: 以下内容对原文作了扩充.}
对每个 \(\beta\), 以 \(\varphi^\beta\) 表示典范投影
\[
  \mathcal L=\bigoplus_\alpha\mathcal O_X^\alpha
    \longrightarrow\mathcal O_X^\beta.
\]
在 \(\mathcal A\) 上的限制.
对 \(x\in X\), 以 \(\mathcal A_x^*\) 表示对偶模
\[
  \operatorname{Hom}_{\mathcal O_{X,x}}
    (\mathcal A_x,\mathcal O_{X,x}).
\]
由于 \(\mathcal A_x\) 是 \(\mathcal L_x\) 的一个直因子,
限制映射
\[
  \operatorname{Hom}_{\mathcal O_{X,x}}
    (\mathcal L_x,\mathcal O_{X,x})
  \longrightarrow
  \operatorname{Hom}_{\mathcal O_{X,x}}
    (\mathcal A_x,\mathcal O_{X,x})
  =\mathcal A_x^*
\]
是满射.
若 \(\mathcal A_x\ne0\), 则由于 \(\mathcal A_x\) 自由且非零,
存在 \(a\in\mathcal A_x\) 及线性形式 \(\varphi\in\mathcal A_x^*\),
使得 \(\varphi(a)=1\).
因此存在 \(\mathcal O_{X,x}\) 中的一族元素 \((\xi^\alpha)\),
使得对每个 \(u\in\mathcal A_{X,x}\), 有
\[
  \varphi(u)=\sum_\alpha\xi^\alpha\varphi^\alpha(u);
\]
这个和是有限的, 因为除有限个 \(\alpha\) 外都有 \(\varphi^\alpha(u)=0\).
将其应用于 \(u=a\), 得到 \(\alpha_1,\ldots,\alpha_n\), 使得
\[
  1=\varphi(a)
   =\xi^{\alpha_1}\varphi^{\alpha_1}(a)
    +\cdots+
    \xi^{\alpha_n}\varphi^{\alpha_n}(a).
\]
于是存在 \(x\) 的一个开邻域 \(U\),
在其上 \(a\) 和各 \(\xi^{\alpha_i}\) 来自截面
\[
  \widetilde a\in\mathcal A(U)
  \qquad\text{及}\qquad
  \widetilde\xi^{\alpha_i}\in\mathcal O_X(U).
\]
等式
\[
  \sum_i
    \widetilde\xi^{\alpha_i}
    \varphi^{\alpha_i}(\widetilde a)=1
  \quad\text{在 }U\text{ 上}
\]
表明对每个 \(y\in U\), 都有 \(\widetilde a_y\ne0\).
(``投射模的支集是开的.'')

\medskip
\phantomsection\label{par:sga3-VIA-2.2}%
\noindent\textbf{2.2.}
记号仍如 \hyperref[par:sga3-VIA-2.1]{2.1}.
显然 \(G'\) 是一个约化 \(k\)-概形.
\hyperref[lem:sga3-VIA-2.1.2]{引理 2.1.2} 表明 \(G'\times_kG'\) 连通,
所以 \((G'\times_kG')_{\mathrm{red}}\) 是 \(G\times_AG\)
中以原点所在的连通分支为底空间的约化子概形.
特别地, 乘法态射
\[
  \mu:G\times_AG\longrightarrow G
\]
诱导态射
\[
  \mu':(G'\times_kG')_{\mathrm{red}}\longrightarrow G'
\]
使 \(G'\) 成为 \((\mathbf{Sch}_{/k})_{\mathrm{red}}\) 中的一个群.

\medskip
\phantomsection\label{par:sga3-VIA-2.2bis}%
\noindent\textbf{2.2.bis.}%
\footnote{\label{note:sga3-VIA-20}编者注: 此段为新增内容,
旨在与第 VI\textsubscript{B} 篇 \S~3 建立联系.}
回顾 (参见第 V 篇), 若 \(P\) 是一个概形,
则 \(\underline P\) 表示 \(P\) 的底拓扑空间.
如下定义 \(G\) 的一个 \(A\)-子函子 \(G^0\):
对每个 \(A\)-概形 \(S\), 置
\[
  G^0(S)
    =\{u\in G(S)\mid u(\underline S)\subset\underline{G'}\}.
\]
设 \(c:G\to G\) 是求逆态射.
由于 \(c(\underline{G'})=\underline{G'}\),
对每个 \(u\in G^0(S)\), 都有 \(c\circ u\in G^0(S)\).
另一方面, 若 \(u,v\in G^0(S)\),
则 \(u\boxtimes v\) 将 \(\underline S\) 映入
\(\underline{G\times_AG}\) 的如下子空间:
两个投影都属于 \(\underline{G'}\) 的那些点所成的子空间.
此子空间等同于 \(G'\times_AG'\) 的底空间,
由\hyperref[lem:sga3-VIA-2.1.2]{引理 2.1.2}, 它是连通的.
所以, \(\mu\circ(u\boxtimes v)\) 将 \(\underline S\)
映入 \(\underline{G'}\).
这证明 \(G^0\) 是 \(G\) 在 \(A\) 上的一个群子函子.

若 \(e\) 所在的连通分支在 \(G\) 中是开的,
则子函子 \(G^0\) 由 \(G\) 在此开子集上诱导的子概形表出,
所以此子概形是 \(G\) 的一个群子概形;
仍将此子概形记作 \(G^0\).
此时, 按 \SGARefLink{prop:sga3-VIB-2.1}{2.1} 的记号,
\[
  G'=(G^0)_{\mathrm{red}},
\]
且拓扑空间 \(\underline{G'}\) 与 \(\underline{G^0}\) 重合.%
\footnote{\label{note:sga3-VIA-21}编者注: 本篇中,
记号 \(G^0\) 仅用于 \(e\) 所在连通分支为开的情形;
第 VI\textsubscript{B} 篇 \S~3 中,
在所有情形下都将以 \(\underline{G^0}\) 表示此连通分支.
这虽是轻微的记号滥用, 但当此连通分支是
\(G\) 的一个开群子概形 \(G^0\) 的底拓扑空间时,
它与以上讨论相容.}

\medskip
\phantomsection\label{par:sga3-VIA-2.3}%
\noindent\textbf{2.3.}
按照 \SGARefLink{par:sga3-VIA-1.1}{1.1} 中的约定,
从现在起再次假设 \(G\) 在 \(A\) 上局部有限型.
则 \(G\) 局部 Noether 型, 因而局部连通;%
\footnote{\label{note:sga3-VIA-22}编者注: 事实上,
Noether 型空间只有有限个连通分支, 因而每个连通分支都是开的;
另见 EGA~I, 6.1.9.}
所以 \(G\) 的每个连通分支都是开的.

以 \(G^0\) 表示 \(G\) 在单位元的连通分支 \(G'\) 上诱导的子概形.
由 \hyperref[par:sga3-VIA-2.2bis]{2.2.bis},
\(G^0\) 是 \(G\) 的一个群子概形,
将其称为 \(G\) 的\emph{单位连通分支 (neutral component)}.
因此, 对每个 \(A\)-概形 \(S\), 有
\[
  G^0(S)
    =\{u\in G(S)\mid
      u(\underline S)\subset\underline{G^0}=\underline{G'}\}.
\]

设 \(G^\alpha\) 是 \(G\) 的任意一个连通分支,
并设
\[
  \nu^\alpha:G^\alpha\times_AG^0\longrightarrow G
\]
是对每个 \(S\in(\mathbf{Sch}_{/A})\), \(g\in G^\alpha(S)\)
及 \(\gamma\in G^0(S)\) 由下式定义的态射:
\[
  \nu^\alpha(S)(g,\gamma)=g\gamma g^{-1}
\]
若 \(e\) 是 \(G\) 的原点,
则 \(\nu^\alpha\) 在 \(G^\alpha\times_A\{e\}\) 上的限制是平凡态射.
由于 \(G^\alpha\times_AG^0\) 由 \hyperref[lem:sga3-VIA-2.1.2]{2.1.2} 可知是连通的,
\(\nu^\alpha\) 通过 \(G^0\) 分解.
所以, 对每个 \(A\)-概形 \(S\),
\(G^0(S)\) 是 \(G(S)\) 的正规子群.
我们证明了以下命题.%
\footnote{\label{note:sga3-VIA-23}编者注:
新增编号 \SGARefLink{prop:sga3-VIA-2.3.1}{2.3.1} 以强调此陈述.
此外, \(G^0\) 甚至是 \(G\) 的特征子群;
见 \SGARefLink{thm:sga3-VIA-2.6.5:item:ii}{2.6.5(ii)}.}

\medskip
\phantomsection\label{prop:sga3-VIA-2.3.1}%
\noindent\textbf{命题 2.3.1.}
\textit{设 \(G\) 是一个局部有限型的 \(A\)-群.
则单位连通分支 \(G^0\) 是 \(G\) 的一个开正规群子概形.}

\medskip
\phantomsection\label{prop:sga3-VIA-2.4}%
\noindent\textbf{命题 2.4.}
\textit{设 \(G\) 是一个局部有限型的 \(A\)-群.}
\begin{enumerate}
\item[\SGAInlineRefTarget{prop:sga3-VIA-2.4:item:i}(i)]
\textit{\(G^0\) 不可约, 且 \(G^0\otimes_Ak\) 在 \(k\) 上几何不可约.}
\item[\SGAInlineRefTarget{prop:sga3-VIA-2.4:item:ii}(ii)]
\textit{\(G^0\) 拟紧, 因而在 \(A\) 上有限型.}
\end{enumerate}

\noindent\footnote{\label{note:sga3-VIA-24}编者注:
陈述中的 ``\(G^0\) 几何不可约'' 已改为
``\(G^0\otimes_Ak\) 在 \(k\) 上几何不可约'', 并扩充了证明.}
\textit{证明.}
\SGARefLink{prop:sga3-VIA-2.4:item:i}{(i)}
由于 \(G^0\) 与 \(G^0\otimes_Ak\) 有相同的底拓扑空间,
只需证明第二个断言.
设 \(\overline k\) 是 \(k\) 的一个代数闭包.
由 \SGARefLink{prop:sga3-VIB-2.2}{2.2},
\((G\otimes_A\overline k)_{\mathrm{red}}\) 是一个局部有限型的约化
\(\overline k\)-群,
因而由 \SGARefLink{prop:sga3-VIA-1.3.1}{1.3.1}, 它在 \(\overline k\) 上平滑.
尤其, \((G\otimes_A\overline k)_{\mathrm{red}}\) 的局部环都是整环.
因此,%
\footnote{\label{note:sga3-VIA-25}编者注: 此处对原文作了扩充,
并新增了对 EGA~I, 6.1.10 的引用.}
由于 \(G\otimes_A\overline k\) 局部 Noether 型,
\(G\otimes_A\overline k\) 的连通分支都不可约
(参见 EGA~I, 6.1.10).
特别地, 连通分支 \(G^0\otimes_A\overline k\) (参见 2.1.1) 不可约.

\SGARefLink{prop:sga3-VIA-2.4:item:ii}{(ii)}
现在证明 \(G^0\) 在 \(A\) 上有限型.
由于 \(G^0\) 在 \(A\) 上局部有限型, 只需证明 \(G^0\) 拟紧.
因 \(G^0\) 不可约, 这由 \SGARefLink{sga3:VIA:0.5.1}{0.5.1} 得到.

\medskip
\phantomsection\label{cor:sga3-VIA-2.4.1}%
\noindent\textbf{推论 2.4.1.}
\textit{\(G\) 的每个连通分支都不可约,%
\footnote{\label{note:sga3-VIA-26}编者注: 须记住,
非单位连通分支一般不是几何连通的.
例如, 若 \(k=\mathbb R\), 则由
\(\mathbb R[X]/(X^3-1)\) 表出的群 \(\mu_{3,\mathbb R}\)
有两个连通分支: \(\{e\}=\operatorname{Spec}\mathbb R\)
和 \(C=\operatorname{Spec}\mathbb R[X]/(X^2+X+1)\);
而 \(C\otimes_{\mathbb R}\mathbb C\) 有两个分支.}
在 \(A\) 上有限型,%
\footnote{\label{note:sga3-VIA-27}编者注: 以下断言为新增内容;
参见第 VI\textsubscript{B} 篇, \SGARefLink{lem:sga3-VIB-1.5}{1.5}.}
且与 \(G^0\) 有相同的维数.}

\medskip
事实上, 可以假设 \(A\) 等于其剩余域 \(k\).
设 \(C\) 是 \(G\) 的一个连通分支,
\(x\) 是 \(C\) 的一个闭点,
\(\kappa(x)\) 是 \(x\) 的剩余域,
\(k'\) 是 \(k\) 的一个包含 \(\kappa(x)\) 的有限正规扩张.
典范投影
\[
  \pi:C\otimes_kk'\longrightarrow C
\]
既开又闭.%
\footnote{\label{note:sga3-VIA-28}编者注: 以下内容对原文作了扩充.}
所以, 若 \(C'\) 是 \(C\otimes_kk'\) 的一个连通分支,
则投影 \(C'\to C\) 是满射.
因而 \(C'\) 包含一个点 \(y\in\pi^{-1}(x)\),
且这样的点在 \(k'\) 上有理 (参见 \hyperref[prop:sga3-VIA-1.2]{1.2} 的证明);
因此 \(C'\) 是 \(G^0\otimes_kk'\) 在右平移 \(r_y\) 下的像.
现在, 由 \hyperref[prop:sga3-VIA-2.4]{2.4},
\(G^0\otimes_kk'\) 在 \(k'\) 上有限型,
且 \(\pi^{-1}(x)\) 是有限集, 基数至多为 \([k':k]\).
所以, \(C\otimes_kk'\) 在 \(k'\) 上有限型,
从而 \(C\) 在 \(k\) 上有限型.

此外, 由于 \(G^0\otimes_kk'\) 由 \SGARefLink{prop:sga3-VIA-2.4}{2.4} 可知不可约,
\(C'\) 也不可约,%
\footnote{\label{note:sga3-VIA-29}编者注: 此处简化了原文.}
而由于投影 \(C'\to C\) 为满射, \(C\) 也不可约.

最后, 以上已知 \(C\otimes_kk'\) 是有限多个
\(G^0\otimes_kk'\) 的平移的不交并.
由于维数在基域扩张下不变
(参见 EGA~IV\textsubscript{2}, 4.1.4),
\(C\) 与 \(G^0\) 的维数相同.
(此外, 由 EGA~IV\textsubscript{2}, 5.2.1,
对每个点 \(g\in G\) 都有 \(\dim_gG=\dim G^0\).)

% Translation coverage: entire assigned component, retaining all source statements and proofs.
% Statistics: source lines=460; target lines=473; inline math=317; display math=31; unresolved translation TODOs=0; independent compilation not performed here.
% Independent prose review: source and translation compared paragraph by paragraph through this component; statements, hypotheses, proofs, notes and diagram captions covered; bibliography reserved for its assigned reviewer; source math and reference keys unchanged.
